\begin{pastexam}[label={px:ust-109-ee-ds-07}]{臺聯大電機類 109 學年度 第 7 題}
（問答；10 分）

(5\%) (a) Sketch the red-black tree by insert 8, 8, 12, 6, 9 into an empty red-black tree.\\
(5\%) (b) Sketch the red-black tree by insert 3, 4, 6, 5, 11, 10 into the red-black tree of part (a).
\end{pastexam}

\begin{pastsoln}
\textbf{題目的問題。} (a) 的 $8$ 出現兩次，(b) 的 $6$ 也在 (a) 插入過，題目沒說重複的鍵怎麼處理。
資結書第~9~章約定重複的鍵忽略；但題目特地寫兩個 $8$，看起來是要插入。 下面的主要解答採用
CLRS 的 RB-Insert 的比較方式：往下找位置時「$k <$ 節點的鍵就往左，否則往右」，所以\textbf{相同
的鍵也插入，放在右子樹}。 第二個 $8$ 記成 $8'$，(b) 的第二個 $6$ 記成 $6'$。 最後一段附上「忽略
重複的鍵」的結果。

插入和修正照資結書第~10.3.4~節（虛擬碼~10.6）：新節點塗紅；父節點是紅的時候看叔叔 $u$——
情況 1（$u$ 紅）重新著色、往上兩層繼續；情況 2（$u$ 黑、$z$ 在內側）對父節點旋轉變成情況 3；
情況 3（$u$ 黑、$z$ 在外側）父節點塗黑、祖父塗紅、對祖父旋轉。 NIL 算黑色。 圖中\textbf{黑底白字
是黑色節點，白底的是紅色節點}。

\textbf{(a)}
\begin{itemize}
\item $8$：空樹，$8$ 當根，塗黑。
\item $8'$：$8' \ge 8$ 往右，父節點 $8$ 是黑的，不用修。
\item $12$：$12 \ge 8$、$12 \ge 8'$，放在 $8'$ 的右邊。 父節點 $8'$ 紅、叔叔（$8$ 的左子）是 NIL，
黑色；$12$ 和 $8'$ 都是右子節點，在外側，情況 3：$8'$ 塗黑、$8$ 塗紅，對 $8$ 左旋。
\item $6$：$6 < 8'$、$6 < 8$，放在 $8$ 的左邊。 父節點 $8$ 紅、叔叔 $12$ 紅，情況 1：$8$、$12$ 塗黑，
$8'$ 塗紅；$8'$ 是根，最後塗回黑色。
\item $9$：$9 \ge 8'$、$9 < 12$，放在 $12$ 的左邊，父節點黑，不用修。
\end{itemize}
\begin{center}
\begin{tikzpicture}[baseline=(current bounding box.north),rn/.style={circle,draw,minimum size=5.6mm,inner sep=0pt,font=\footnotesize},bn/.style={rn,fill=black!85,text=white,font=\footnotesize\bfseries}]
\node[bn] (v0) at (0.00,-0.00) {8};
\node[font=\small,anchor=north] at (0.00,-0.45) {插入 8 後};
\end{tikzpicture}
\hspace{0.8cm}
\begin{tikzpicture}[baseline=(current bounding box.north),rn/.style={circle,draw,minimum size=5.6mm,inner sep=0pt,font=\footnotesize},bn/.style={rn,fill=black!85,text=white,font=\footnotesize\bfseries}]
\node[bn] (v0) at (0.00,-0.00) {8};
\node[rn] (v1) at (0.62,-0.85) {$8'$};
\draw (v0)--(v1);
\node[font=\small,anchor=north] at (0.31,-1.30) {插入 $8'$ 後};
\end{tikzpicture}
\hspace{0.8cm}
\begin{tikzpicture}[baseline=(current bounding box.north),rn/.style={circle,draw,minimum size=5.6mm,inner sep=0pt,font=\footnotesize},bn/.style={rn,fill=black!85,text=white,font=\footnotesize\bfseries}]
\node[rn] (v0) at (0.00,-0.85) {8};
\node[bn] (v1) at (0.62,-0.00) {$8'$};
\node[rn] (v2) at (1.24,-0.85) {12};
\draw (v1)--(v0);
\draw (v1)--(v2);
\node[font=\small,anchor=north] at (0.62,-1.30) {插入 12 後};
\end{tikzpicture}
\end{center}
\begin{center}
\begin{tikzpicture}[baseline=(current bounding box.north),rn/.style={circle,draw,minimum size=5.6mm,inner sep=0pt,font=\footnotesize},bn/.style={rn,fill=black!85,text=white,font=\footnotesize\bfseries}]
\node[rn] (v3) at (0.00,-1.70) {6};
\node[bn] (v0) at (0.62,-0.85) {8};
\node[bn] (v1) at (1.24,-0.00) {$8'$};
\node[bn] (v2) at (1.86,-0.85) {12};
\draw (v0)--(v3);
\draw (v1)--(v0);
\draw (v1)--(v2);
\node[font=\small,anchor=north] at (0.93,-2.15) {插入 6 後};
\end{tikzpicture}
\hspace{0.8cm}
\begin{tikzpicture}[baseline=(current bounding box.north),rn/.style={circle,draw,minimum size=5.6mm,inner sep=0pt,font=\footnotesize},bn/.style={rn,fill=black!85,text=white,font=\footnotesize\bfseries}]
\node[rn] (v3) at (0.00,-1.70) {6};
\node[bn] (v0) at (0.62,-0.85) {8};
\node[bn] (v1) at (1.24,-0.00) {$8'$};
\node[rn] (v4) at (1.86,-1.70) {9};
\node[bn] (v2) at (2.48,-0.85) {12};
\draw (v0)--(v3);
\draw (v1)--(v0);
\draw (v1)--(v2);
\draw (v2)--(v4);
\node[font=\small,anchor=north] at (1.24,-2.15) {插入 9 後};
\end{tikzpicture}
\end{center}
(a) 的答案是插入 $9$ 後的樹：根 $8'$（黑），左子 $8$（黑，左子 $6$ 紅），右子 $12$（黑，左子 $9$ 紅）。
左旋之後 $8$ 跑到 $8'$ 的左邊，相同的鍵出現在左子樹；有重複鍵的二元搜尋樹本來就只能保證
「左 $\le$ 節點 $\le$ 右」，搜尋時要注意。

\textbf{(b)}
\begin{itemize}
\item $3$：一路往左，放在 $6$ 的左邊。 父節點 $6$ 紅、祖父 $8$、叔叔（$8$ 的右子）NIL 黑；$3$、$6$
都是左子節點，情況 3：$6$ 塗黑、$8$ 塗紅，對 $8$ 右旋，這一段變成 $6$（黑）下面掛 $3$、$8$（紅）。
\item $4$：放在 $3$ 的右邊。 父節點 $3$ 紅、叔叔 $8$ 紅，情況 1：$3$、$8$ 塗黑，$6$ 塗紅。 $6$ 的
父節點 $8'$ 是黑的，停止。
\item $6'$：$6' < 8'$ 往左到 $6$，$6' \ge 6$ 往右到 $8$，$6' < 8$ 放在 $8$ 的左邊。 父節點黑，不用修。
\item $5$：$5 < 8'$、$5 < 6$、$5 \ge 3$、$5 \ge 4$，放在 $4$ 的右邊。 父節點 $4$ 紅、祖父 $3$、叔叔
（$3$ 的左子）NIL 黑；$5$、$4$ 都是右子節點，情況 3（左右對調）：$4$ 塗黑、$3$ 塗紅，對 $3$ 左旋。
\item $11$：放在 $9$ 的右邊。 父節點 $9$ 紅、祖父 $12$、叔叔 NIL 黑；$9$ 是左子、$11$ 是右子，在
內側，情況 2：對 $9$ 左旋，$9$ 變成 $11$ 的左子；接著情況 3：$11$ 塗黑、$12$ 塗紅，對 $12$ 右旋。
\item $10$：$10 \ge 8'$、$10 < 11$、$10 \ge 9$，放在 $9$ 的右邊。 父節點 $9$ 紅、叔叔 $12$ 紅，
情況 1：$9$、$12$ 塗黑，$11$ 塗紅。 $11$ 的父節點 $8'$ 黑，停止。
\end{itemize}
\begin{center}
\begin{tikzpicture}[baseline=(current bounding box.north),rn/.style={circle,draw,minimum size=5.6mm,inner sep=0pt,font=\footnotesize},bn/.style={rn,fill=black!85,text=white,font=\footnotesize\bfseries}]
\node[rn] (v5) at (0.00,-1.70) {3};
\node[bn] (v3) at (0.62,-0.85) {6};
\node[rn] (v0) at (1.24,-1.70) {8};
\node[bn] (v1) at (1.86,-0.00) {$8'$};
\node[rn] (v4) at (2.48,-1.70) {9};
\node[bn] (v2) at (3.10,-0.85) {12};
\draw (v3)--(v5);
\draw (v3)--(v0);
\draw (v1)--(v3);
\draw (v1)--(v2);
\draw (v2)--(v4);
\node[font=\small,anchor=north] at (1.55,-2.15) {插入 3 後};
\end{tikzpicture}
\hspace{0.8cm}
\begin{tikzpicture}[baseline=(current bounding box.north),rn/.style={circle,draw,minimum size=5.6mm,inner sep=0pt,font=\footnotesize},bn/.style={rn,fill=black!85,text=white,font=\footnotesize\bfseries}]
\node[bn] (v5) at (0.00,-1.70) {3};
\node[rn] (v6) at (0.62,-2.55) {4};
\node[rn] (v3) at (1.24,-0.85) {6};
\node[bn] (v0) at (1.86,-1.70) {8};
\node[bn] (v1) at (2.48,-0.00) {$8'$};
\node[rn] (v4) at (3.10,-1.70) {9};
\node[bn] (v2) at (3.72,-0.85) {12};
\draw (v5)--(v6);
\draw (v3)--(v5);
\draw (v3)--(v0);
\draw (v1)--(v3);
\draw (v1)--(v2);
\draw (v2)--(v4);
\node[font=\small,anchor=north] at (1.86,-3.00) {插入 4 後};
\end{tikzpicture}
\end{center}
\begin{center}
\begin{tikzpicture}[baseline=(current bounding box.north),rn/.style={circle,draw,minimum size=5.6mm,inner sep=0pt,font=\footnotesize},bn/.style={rn,fill=black!85,text=white,font=\footnotesize\bfseries}]
\node[bn] (v5) at (0.00,-1.70) {3};
\node[rn] (v6) at (0.62,-2.55) {4};
\node[rn] (v3) at (1.24,-0.85) {6};
\node[rn] (v7) at (1.86,-2.55) {$6'$};
\node[bn] (v0) at (2.48,-1.70) {8};
\node[bn] (v1) at (3.10,-0.00) {$8'$};
\node[rn] (v4) at (3.72,-1.70) {9};
\node[bn] (v2) at (4.34,-0.85) {12};
\draw (v5)--(v6);
\draw (v3)--(v5);
\draw (v3)--(v0);
\draw (v0)--(v7);
\draw (v1)--(v3);
\draw (v1)--(v2);
\draw (v2)--(v4);
\node[font=\small,anchor=north] at (2.17,-3.00) {插入 $6'$ 後};
\end{tikzpicture}
\hspace{0.8cm}
\begin{tikzpicture}[baseline=(current bounding box.north),rn/.style={circle,draw,minimum size=5.6mm,inner sep=0pt,font=\footnotesize},bn/.style={rn,fill=black!85,text=white,font=\footnotesize\bfseries}]
\node[rn] (v5) at (0.00,-2.55) {3};
\node[bn] (v6) at (0.62,-1.70) {4};
\node[rn] (v8) at (1.24,-2.55) {5};
\node[rn] (v3) at (1.86,-0.85) {6};
\node[rn] (v7) at (2.48,-2.55) {$6'$};
\node[bn] (v0) at (3.10,-1.70) {8};
\node[bn] (v1) at (3.72,-0.00) {$8'$};
\node[rn] (v4) at (4.34,-1.70) {9};
\node[bn] (v2) at (4.96,-0.85) {12};
\draw (v6)--(v5);
\draw (v6)--(v8);
\draw (v3)--(v6);
\draw (v3)--(v0);
\draw (v0)--(v7);
\draw (v1)--(v3);
\draw (v1)--(v2);
\draw (v2)--(v4);
\node[font=\small,anchor=north] at (2.48,-3.00) {插入 5 後};
\end{tikzpicture}
\end{center}
\begin{center}
\begin{tikzpicture}[baseline=(current bounding box.north),rn/.style={circle,draw,minimum size=5.6mm,inner sep=0pt,font=\footnotesize},bn/.style={rn,fill=black!85,text=white,font=\footnotesize\bfseries}]
\node[rn] (v5) at (0.00,-2.55) {3};
\node[bn] (v6) at (0.62,-1.70) {4};
\node[rn] (v8) at (1.24,-2.55) {5};
\node[rn] (v3) at (1.86,-0.85) {6};
\node[rn] (v7) at (2.48,-2.55) {$6'$};
\node[bn] (v0) at (3.10,-1.70) {8};
\node[bn] (v1) at (3.72,-0.00) {$8'$};
\node[rn] (v4) at (4.34,-1.70) {9};
\node[bn] (v9) at (4.96,-0.85) {11};
\node[rn] (v2) at (5.58,-1.70) {12};
\draw (v6)--(v5);
\draw (v6)--(v8);
\draw (v3)--(v6);
\draw (v3)--(v0);
\draw (v0)--(v7);
\draw (v1)--(v3);
\draw (v1)--(v9);
\draw (v9)--(v4);
\draw (v9)--(v2);
\node[font=\small,anchor=north] at (2.79,-3.00) {插入 11 後};
\end{tikzpicture}
\hspace{0.8cm}
\begin{tikzpicture}[baseline=(current bounding box.north),rn/.style={circle,draw,minimum size=5.6mm,inner sep=0pt,font=\footnotesize},bn/.style={rn,fill=black!85,text=white,font=\footnotesize\bfseries}]
\node[rn] (v5) at (0.00,-2.55) {3};
\node[bn] (v6) at (0.62,-1.70) {4};
\node[rn] (v8) at (1.24,-2.55) {5};
\node[rn] (v3) at (1.86,-0.85) {6};
\node[rn] (v7) at (2.48,-2.55) {$6'$};
\node[bn] (v0) at (3.10,-1.70) {8};
\node[bn] (v1) at (3.72,-0.00) {$8'$};
\node[bn] (v4) at (4.34,-1.70) {9};
\node[rn] (v10) at (4.96,-2.55) {10};
\node[rn] (v9) at (5.58,-0.85) {11};
\node[bn] (v2) at (6.20,-1.70) {12};
\draw (v6)--(v5);
\draw (v6)--(v8);
\draw (v3)--(v6);
\draw (v3)--(v0);
\draw (v0)--(v7);
\draw (v1)--(v3);
\draw (v1)--(v9);
\draw (v4)--(v10);
\draw (v9)--(v4);
\draw (v9)--(v2);
\node[font=\small,anchor=north] at (3.10,-3.00) {插入 10 後};
\end{tikzpicture}
\end{center}
(b) 的答案是插入 $10$ 後的樹。 檢查：根黑；沒有紅節點的子節點是紅的；從根到任何一個 NIL 都經過
$2$ 個黑色節點（不算 NIL），例如 $8' \to 6 \to 4 \to 3$ 是 $8'$、$4$，$8' \to 11 \to 12$ 是 $8'$、$12$。

\textbf{如果重複的鍵忽略不插入}（照資結書第~9~章的約定），(a) 等於依序插入 $8, 12, 6, 9$，
(b) 再插入 $3, 4, 5, 11, 10$：插入 $9$ 時叔叔 $6$ 紅，情況 1；插入 $4$ 是情況 2 加情況 3（對 $3$
左旋、對 $6$ 右旋）；插入 $5$ 時叔叔 $3$ 紅，情況 1；$11$、$10$ 和上面一樣。 結果：
\begin{center}
\begin{tikzpicture}[baseline=(current bounding box.north),rn/.style={circle,draw,minimum size=5.6mm,inner sep=0pt,font=\footnotesize},bn/.style={rn,fill=black!85,text=white,font=\footnotesize\bfseries}]
\node[bn] (v3) at (0.00,-0.85) {6};
\node[bn] (v0) at (0.62,-0.00) {8};
\node[rn] (v4) at (1.24,-1.70) {9};
\node[bn] (v2) at (1.86,-0.85) {12};
\draw (v0)--(v3);
\draw (v0)--(v2);
\draw (v2)--(v4);
\node[font=\small,anchor=north] at (0.93,-2.15) {(a) 的結果};
\end{tikzpicture}
\hspace{1cm}
\begin{tikzpicture}[baseline=(current bounding box.north),rn/.style={circle,draw,minimum size=5.6mm,inner sep=0pt,font=\footnotesize},bn/.style={rn,fill=black!85,text=white,font=\footnotesize\bfseries}]
\node[bn] (v5) at (0.00,-1.70) {3};
\node[rn] (v6) at (0.62,-0.85) {4};
\node[rn] (v8) at (1.24,-2.55) {5};
\node[bn] (v3) at (1.86,-1.70) {6};
\node[bn] (v0) at (2.48,-0.00) {8};
\node[bn] (v4) at (3.10,-1.70) {9};
\node[rn] (v10) at (3.72,-2.55) {10};
\node[rn] (v9) at (4.34,-0.85) {11};
\node[bn] (v2) at (4.96,-1.70) {12};
\draw (v6)--(v5);
\draw (v6)--(v3);
\draw (v3)--(v8);
\draw (v0)--(v6);
\draw (v0)--(v9);
\draw (v4)--(v10);
\draw (v9)--(v4);
\draw (v9)--(v2);
\node[font=\small,anchor=north] at (2.48,-3.00) {(b) 的結果};
\end{tikzpicture}
\end{center}
作答時把採用的處理方式寫在最前面，閱卷才知道你的樹怎麼來的。
\end{pastsoln}
